\section{In cammino verso la Santa Cresima}

Da sabato 14 febbraio a lunedì 16 febbraio, 74 ragazzi di seconda media, di Maserada, Varago e Candelù, insieme  a Don Federico, a noi catechiste  e ad alcuni accompagnatori abbiamo partecipato ad un'uscita di tre giorni ad Assisi.  L'idea ha preso forma all'inizio dell'anno catechistico ed è stata subito accolta con entusiasmo. Tutti insieme  ci siamo quindi  messi al lavoro per organizzare al meglio l'esperienza e recuperare fondi per contenere le spese. A fine dicembre grazie al prezioso contributo dell'Associazione ciclistica "Ottavio Zuliani" che ci ha gentilmente messo a disposizione un pulmino da 9 posti, ci siamo recati ad Assisi per pianificare  l'itinerario e le varie attività. Per cercare di ammortizzare i costi dell'uscita, nelle tre parrocchie, le famiglie dei ragazzi  si sono messe all'opera con una vendita di dolci e con una serata dimostrativa di materassi: anche a loro va il nostro ringraziamento per il supporto e il sostegno.

E così sulle orme di S.Francesco, S.Chiara e S.Carlo Acutis  il nostro cammino ha iniziato a prendere forma: tre strade diverse di santità , tre santi che si sono lasciati plasmare dall'amore di Gesù e lo hanno testimoniato con la  vita.  
Pronti, partenza, via....

Sabato mattina alle ore 6.00 siamo partiti con 2 corriere. La pioggia ci ha accompagnati durante quasi tutto il viaggio di andata, ma in autobus c'era il sole: tra ragazzi e adulti si respirava un'aria gioiosa e allegra e così,  tra canti e chiacchiere, siamo arrivati a Rivotorto. Nel Santuario del Sacro Tugurio, alle porte di Assisi, abbiamo iniziato la nostra attività  con una scenetta che ci ha introdotti nella vita di S. Francesco. Nel primo pomeriggio poi abbiamo visitato la Basilica di S. Maria degli Angeli e ascoltato la testimonianza di  Fra Luca che ha saputo catturare la nostra attenzione e sollecitato gli animi di tutti noi. All'uscita il cielo era sereno, perciò ci siamo incamminati verso Assisi percorrendo la Strada Mattonata. Nella serata di sabato, abbiamo poi vissuto un altro momento particolarmente intenso e toccante con la  visita a porte chiuse della Basilica e della tomba di S. Francesco. Fra' Nicola ci ha accompagnati in un suggestivo viaggio tra gli affreschi della basilica superiore e in un intenso momento di raccoglimento davanti  alla tomba del Santo.

Domenica mattina, Don Federico  e Don  Graziano hanno celebrato l'Eucarestia nella Basilica Inferiore; un momento di grazia per i ragazzi cresimandi di  Maserada e  Postioma; una bella occasione di fraternità, un segno che ci ha fatti sentire parte di una grande famiglia. La tappa successiva è stata al Santuario della Spogliazione dove è custodito il corpo di S. Carlo Acutis; un ragazzo straordinario che ha saputo vivere la sua Santità nella quotidianità. Nel pomeriggio poi visita al Santuario di S.Damiano dove ci accolti Frate Antonio e successivamente alla basilica di S.Chiara con una preghiera davanti al crocifisso di San Damiano. Per la serata, ci aspettava un'altra sfida:  muniti di torce e ben vestiti, ci siamo incamminati verso la Rocca dove avevamo appuntamento con i ragazzi di Postioma per fare insieme un grande gioco notturno e un momento di preghiera che si è concluso illuminando il cielo con una croce formata da tante piccole luci. 

\figurawrap{0.65\textwidth}{cresime}

Lunedì mattina: partenza da Assisi verso La Verna. Qui, la giornata fredda, umida e nebbiosa ci ha aiutati a ritrovare un po' noi stessi.  Grazie alla guida di Don Federico  ci siamo dati del tempo per la  meditazione personale, per ascoltare il nostro cuore e raccogliere i frutti di questa esperienza. Dopo un momento di condivisione e un pranzo caldo, con la gioia nel cuore e gli occhi che brillavano più del solito,  ci siamo incamminati verso le corriere per il viaggio di ritorno.

Per noi catechiste è stata un'esperienza intensa e ricca di vita. Un grazie particolare a chi da casa ha contribuito in vario modo per darci la possibilità di vivere questa meravigliosa avventura e alle tante persone chi ci hanno sostenuto e accompagnato con la preghiera: abbiamo sperimentato che non siamo soli.  E' stato un tempo ricco di grazia, di tanti segni piccoli e grandi che abbiamo accolto con occhi e cuore aperti, nulla è stato casuale. Unendo le nostre forze abbiamo sperimentato la ricchezza della collaborazione tra parrocchie e siamo riusciti a formare una bella  squadra.  Abbiamo vissuto tre giorni intensi, pieni  di emozioni, di risate, di fatiche, di spiritualità e di pace.  Ora carichi e felici, custodendo nel cuore tanti bei ricordi di questa bella esperienza, ci stiamo preparando per ricevere domenica 17 maggio il Dono più grande: Lo Spirito Santo. 

\firma{Le catechiste di seconda media}
